% Part: sets-functions-relations
% Chapter: relations
% Section: orders

\documentclass[../../../include/open-logic-section]{subfiles}

\begin{document}

\olfileid[fr]{sfr}{rel}{ord}
\olsection{Ordres}

\begin{explain}
Dans bien des comparaisons, nous décrivons certains objets comme
« inférieurs », « égaux » ou « supérieurs » à d’autres, sous un certain
rapport. Ces comparaisons font intervenir des relations d’\emph{ordre}.
Il en existe cependant plusieurs types. Certaines exigent que deux objets
quelconques soient comparables, d’autres non. Certaines comprennent
l’identité (comme~$\le$), d’autres l’excluent (comme~$<$).
Il nous sera utile de les classer.
\end{explain}

\begin{defn}[Préordre]
Une relation à la fois réflexive et transitive est appelée un
\emph{préordre}.
\end{defn}

\begin{defn}[Ordre partiel]
Un préordre qui est aussi antisymétrique est appelé un \emph{ordre partiel}.
\end{defn}

\begin{defn}[Ordre linéaire]\ollabel{def:linearorder}
Un ordre partiel qui est aussi connexe est appelé un \emph{ordre total}
ou \emph{ordre linéaire}.
\end{defn}

\begin{ex}
Tout ordre linéaire est aussi un ordre partiel, et tout ordre partiel
est aussi un préordre, mais les réciproques sont fausses.
La relation universelle sur~$A$ est un préordre, car elle est réflexive
et transitive. Mais, si $A$ possède plus d’un !!{element}, cette relation
n’est pas antisymétrique ; ce n’est donc pas un ordre partiel.
\end{ex}

\begin{ex}
Considérons sur~$\Bin^*$ la relation $\preccurlyeq$, qui se lit
\emph{est de longueur inférieure ou égale à} : $x \preccurlyeq y$ si et
seulement si $\len{x} \le \len{y}$. C’est un préordre (réflexif et
transitif), qui est même connexe, mais ce n’est pas un ordre partiel,
car il n’est pas antisymétrique. Ainsi, $01 \preccurlyeq 10$ et
$10 \preccurlyeq 01$, alors que $01 \neq 10$.
\end{ex}

\begin{ex}
La relation $\subseteq$ sur un ensemble d’ensembles est un ordre partiel
important. Ce n’est pas en général un ordre linéaire : si $a \neq b$ et
que l’on considère $\Pow{\{a, b\}} = \{\emptyset, \{a\}, \{b\}, \{a,b\}\}$,
on constate que $\{a\} \nsubseteq \{b\}$, que $\{a\} \neq \{b\}$ et que
$\{b\} \nsubseteq \{a\}$.
\end{ex}

\begin{ex}
La relation de \emph{divisibilité} fournit un ordre partiel qui n’est pas
linéaire. Pour des entiers $n$ et~$m$, on écrit $n \mid m$ pour dire que
$n$ divise $m$ sans reste, c’est-à-dire qu’il existe un entier~$k$ tel que
$m = kn$. Sur~$\Nat$, c’est un ordre partiel, mais pas un ordre linéaire :
par exemple, $2 \nmid 3$ et $3 \nmid 2$. Sur $\Int$, la divisibilité
n’est qu’un préordre, car elle n’est pas antisymétrique : $1 \mid -1$ et
$-1 \mid 1$, alors que $1 \neq -1$.
\end{ex}

\begin{ex}
La relation de \emph{prolongement} sur l’ensemble de suites~$A^*$ se
définit ainsi : $s \sqsubseteq s'$ si et seulement si $s = \emptyseq$
(la suite vide), $s = s'$, ou bien $s = \tuple{s_1, \dots, s_n}$ et
$s' = \tuple{s_1, \dots, s_n, s_{n+1}, \dots, s_m}$.
Lorsque $s \sqsubseteq s'$, on dit aussi que $s$ est un \emph{segment initial}
de~$s'$. La relation de prolongement sur $A^*$ est un ordre partiel.
Elle n’est pas linéaire dès que l’alphabet contient deux symboles distincts :
si $a \neq b$, alors $ab \not\sqsubseteq ba$ et $ba \not\sqsubseteq ab$.
\footnote{Précisions éditoriales : avec au plus un symbole, cet ordre est
linéaire. On considère ici de véritables suites finies, dont la longueur
fait partie de l’identité, conformément à la précision sur le codage des
mots donnée au chapitre précédent.}
\end{ex}

\begin{defn}[Ordre strict]
Un \emph{ordre strict} est une relation irréflexive, asymétrique et transitive.
\end{defn}

\begin{defn}[Ordre linéaire strict]\ollabel{def:strictlinearorder}
Un ordre strict qui est aussi connexe est appelé un \emph{ordre total strict}
ou \emph{ordre linéaire strict}.
\end{defn}

\begin{ex}
$\le$ est l’ordre linéaire correspondant à l’ordre linéaire strict~$<$.
$\subseteq$ est l’ordre partiel correspondant à l’ordre strict~$\subsetneq$.
\end{ex}

On peut transformer tout ordre strict $R$ sur~$A$ en un ordre partiel
en lui ajoutant la diagonale $\Id{A}$, c’est-à-dire tous les
couples~$\tuple{x, x}$. On obtient ainsi la \emph{clôture réflexive} de~$R$.
Réciproquement, on obtient un ordre strict à partir d’un ordre partiel
en en retirant~$\Id{A}$. Les deux résultats suivants le précisent.

\begin{prop}\ollabel{prop:stricttopartial}
Si $R$ est un ordre strict sur~$A$, alors $R^+ = R \cup \Id{A}$ est un
ordre partiel. De plus, si $R$ est un ordre linéaire strict, $R^+$ est
un ordre linéaire.
\end{prop}

\begin{proof}
Supposons que $R$ soit un ordre strict : $R \subseteq A^2$ et $R$ est
irréflexive, asymétrique et transitive. Posons $R^+ = R \cup \Id{A}$.
Il faut montrer que $R^+$ est réflexive, antisymétrique et transitive.

$R^+$ est manifestement réflexive, puisque
$\tuple{x, x} \in \Id{A} \subseteq R^+$ pour tout $x \in A$.

Pour établir l’antisymétrie de $R^+$, supposons, par l’absurde, que
$R^+xy$ et $R^+yx$, mais que $x \neq y$. Puisque
$\tuple{x,y} \in R \cup \Id{A}$ et $\tuple{x, y} \notin \Id{A}$,
on a nécessairement $\tuple{x, y} \in R$, c’est-à-dire $Rxy$.
De même, on a~$Ryx$. Cela contredit l’hypothèse selon laquelle $R$
est asymétrique.

Pour la transitivité, supposons $R^+xy$ et $R^+yz$. Si l’on a à la fois
$\tuple{x, y} \in R$ et $\tuple{y,z} \in R$, alors
$\tuple{x, z} \in R$, puisque $R$ est transitive. Sinon,
$\tuple{x, y} \in \Id{A}$, c’est-à-dire $x = y$, ou bien
$\tuple{y, z} \in \Id{A}$, c’est-à-dire $y = z$.
Dans le premier cas, on a $R^+yz$ par hypothèse et $x = y$, donc $R^+xz$.
Le second cas se traite de la même manière. Dans tous les cas, $R^+xz$ :
$R^+$ est donc aussi transitive.

Pour la dernière affirmation, supposons que $R$ soit en outre connexe.
Pour tous $x \neq y$, on a donc $Rxy$ ou~$Ryx$, c’est-à-dire
$\tuple{x, y} \in R$ ou $\tuple{y, x} \in R$.
Puisque $R \subseteq R^+$, cela reste vrai pour $R^+$, qui est donc
connexe elle aussi.
\end{proof}

\begin{prop}\ollabel{prop:partialtostrict}
Si $R$ est un ordre partiel sur~$A$, alors $R^- = R \setminus \Id{A}$
est un ordre strict. De plus, si $R$ est un ordre linéaire, $R^-$ est
un ordre linéaire strict.
\end{prop}

\begin{proof}
La démonstration est laissée en exercice.
\end{proof}

\begin{prob}
Démontrez la \olref[sfr][rel][ord]{prop:partialtostrict}.
\end{prob}

Le résultat élémentaire suivant établit que les ordres linéaires stricts
possèdent une propriété analogue à l’extensionnalité.

\begin{prop}\ollabel{prop:extensionality-strictlinearorders}
Si $<$ est un ordre linéaire strict sur $A$, alors :
\[
  (\forall a, b \in A)((\forall x \in A)(x < a \liff x < b) \lif a = b).
\]
\end{prop}

\begin{proof}
Supposons $(\forall x \in A)(x < a \liff x < b)$. Si $a < b$, alors
$a < a$, ce qui contredit l’irréflexivité de $<$ ; donc $a \nless b$.
De la même manière, $b \nless a$. Puisque $<$ est connexe, on a $a = b$.
\end{proof}

\end{document}
